% BAGEL 两卡实测结果：全部结果矩阵、最优配置对比、等量比较、首请求
\setpart{BAGEL 实测 · 两卡结果}

\begin{frame}[t]{两卡全部实测结果：卡 1、2，稳态中位数，秒}
\scriptsize\def\codesize{\relax}
\renewcommand{\arraystretch}{1.12}\setlength{\tabcolsep}{3pt}
\begin{tabularx}{\textwidth}{>{\bfseries}p{5.9cm} p{0.9cm} Y Y Y Y}
\toprule
\rowcolor{Mist}
配置 & 框架 & T2I 512²·50 步 & I2T 1 token & I2T 64 token & I2I 512²·50 步 \\
\midrule
M* CFG 并行 · compile 开 · 缓存开 · 无 graph & \textcolor{Teal}{M*} & \cellcolor{GoldBg}\textbf{\textcolor{Teal}{7.68}}\,$\star$ & 0.085 & 2.46 & \cellcolor{GoldBg}\textbf{\textcolor{Teal}{16.79}}\,$\star$ \\
M* CFG 并行 · compile 关（诊断）· 缓存开 & \textcolor{Teal}{M*} & 8.48 & \cellcolor{GoldBg}\textbf{\textcolor{Teal}{0.083}}\,$\star$ & \cellcolor{GoldBg}\textbf{\textcolor{Teal}{2.17}}\,$\star$ & 18.25 \\
M* TP2 · compile 开 · 缓存开 · 无 graph & \textcolor{Teal}{M*} & 12.96 & 0.330 & 3.16 & 21.21 \\
M* CFG 并行 · compile 开 · 缓存关 & \textcolor{Teal}{M*} & -- & 0.293 & 2.66 & \cellcolor{GoldBg}\textbf{\textcolor{Teal}{16.82}}\,$\star$ \\
\midrule
Omni 单 stage CFG 并行 · 无 compile & \textcolor{Orange}{Omni} & \cellcolor{GoldBg}\textbf{\textcolor{Teal}{9.98}}\,$\star$ & -- & -- & -- \\
Omni 单 stage TP2 · 无 graph & \textcolor{Orange}{Omni} & 17.04 & -- & -- & \cellcolor{GoldBg}\textbf{\textcolor{Teal}{28.08}}\,$\star$ \\
Omni two-stage · 缓存关 · eager & \textcolor{Orange}{Omni} & 17.99 & 0.193 & 1.88 & \textcolor{Muted}{30.76\textsuperscript{†}} \\
Omni two-stage · 缓存关 · Thinker graph & \textcolor{Orange}{Omni} & -- & 0.191 & 1.86 & -- \\
Omni two-stage · 缓存开 · eager & \textcolor{Orange}{Omni} & -- & \cellcolor{GoldBg}\textbf{\textcolor{Teal}{0.053}}\,$\star$ & \cellcolor{GoldBg}\textbf{\textcolor{Teal}{1.74}}\,$\star$ & -- \\
\bottomrule
\end{tabularx}

\vspace{3pt}
\begin{itemize}\scriptsize\setlength\itemsep{0pt}
  \item 配置列即该行的完整定义，见第 3 部分。$\star$ = 同一框架内该 case 的最佳（1\% 以内并列）。n = 5，min--max 见附录 A。
  \item 图生文两边输入规模一致（M* 1369 patch，Omni 1386 token）；图像编辑两边指令相同。--：未测或不支持；\textsuperscript{†}灰图，只保留耗时。
\end{itemize}
\takeaway{M* 最优：文生图 \val{final/mstar-cfgpar/t2i} s，图像编辑 \val{final/mstar-cfgpar/i2i} s。vLLM-Omni 最优：文生图 \val{final/omni-single-cfgpar2/t2i} s，图生文 \val{fact/omni-two-stage-prefixcache/i2t1} s（1 token）/ \val{fact/omni-two-stage-prefixcache/i2t64} s（64 token，缓存开）。}
\end{frame}

\newcommand{\barpair}[6]{% y, label, mstar value, omni value, mstar tag, omni tag
    \pgfmathsetmacro{\mx}{max(#3,#4)}
    \pgfmathsetmacro{\wm}{#3/\mx*6.6}
    \pgfmathsetmacro{\wo}{#4/\mx*6.6}
    \node[anchor=east,text width=3.6cm,align=right,font=\scriptsize\bfseries] at (3.7,#1-0.34) {#2};
    \fill[Teal] (3.9,#1-0.04) rectangle ++(\wm,-0.28);
    \node[anchor=west,font=\tiny] at (3.95+\wm,#1-0.18) {#3\,s\ \textcolor{Muted}{#5}};
    \fill[Orange] (3.9,#1-0.36) rectangle ++(\wo,-0.28);
    \node[anchor=west,font=\tiny] at (3.95+\wo,#1-0.50) {#4\,s\ \textcolor{Muted}{#6}};}

\begin{frame}[t]{各 case 的最优配置对比（柱长按 case 内较大值归一）}
\begin{tikzpicture}[x=1cm,y=.72cm,font=\scriptsize]
  \fill[Teal] (3.9,.85) rectangle ++(.3,-.25); \node[anchor=west,font=\scriptsize] at (4.3,.72) {M*（CFG 并行，compile 开）};
  \fill[Orange] (8.2,.85) rectangle ++(.3,-.25); \node[anchor=west,font=\scriptsize] at (8.6,.72) {vLLM-Omni};
  \barpair{0}{文生图 T2I}{\raw{final/mstar-cfgpar/t2i}}{\raw{final/omni-single-cfgpar2/t2i}}{CFG 并行}{单 stage CFG 并行}
  \barpair{-1.0}{图像编辑 I2I}{\raw{final/mstar-cfgpar/i2i}}{\raw{final/omni-single-tp2/i2i}}{CFG 并行}{单 stage TP2}
  \barpair{-2.0}{图生文 1 token，缓存关}{\raw{fact/mstar-cfgpar-nocache/i2t1}}{\raw{final/omni-two-stage/i2t1}}{CFG 并行}{two-stage，eager}
  \barpair{-3.0}{图生文 1 token，缓存开}{\raw{final/mstar-cfgpar/i2t1}}{\raw{fact/omni-two-stage-prefixcache/i2t1}}{}{two-stage + 前缀缓存}
  \barpair{-4.0}{图生文 64 token，缓存关}{\raw{fact/mstar-cfgpar-nocache/i2t64}}{\raw{final/omni-two-stage/i2t64}}{CFG 并行}{two-stage，eager}
  \barpair{-5.0}{图生文 64 token，缓存开}{\raw{final/mstar-cfgpar/i2t64}}{\raw{fact/omni-two-stage-prefixcache/i2t64}}{}{two-stage + 前缀缓存}
\end{tikzpicture}
\takeaway{vLLM-Omni 相对 M* 的耗时：文生图为 \pct{bestTti}，图像编辑为 \pct{bestIti}。图生文 1 token 为 \pct{i2tOff1}（缓存关）、\pct{i2tOn1}（缓存开）；64 token 为 \pct{i2tOff64}、\pct{i2tOn64}。}
\end{frame}

\begin{frame}[t]{图生文：输入与缓存设置对齐后的结果}
\scriptsize\def\codesize{\relax}
\renewcommand{\arraystretch}{1.12}\setlength{\tabcolsep}{3pt}
\begin{tabularx}{\textwidth}{>{\bfseries}p{4.6cm} Y Y Y Y}
\toprule
\rowcolor{Mist}
M* 配置 & 1 token，4900 patch（早期） & 1 token，1369 patch（本次） & 64 token，4900 patch（早期） & 64 token，1369 patch（本次） \\
\midrule
CFG 并行（compile 开） & \textcolor{Muted}{0.193} & \textbf{\textcolor{Teal}{0.085}} & \textcolor{Muted}{2.55} & \textbf{\textcolor{Teal}{2.46}} \\
CFG 并行（compile 关） & \textcolor{Muted}{0.200} & \textbf{\textcolor{Teal}{0.083}} & \textcolor{Muted}{2.29} & \textbf{\textcolor{Teal}{2.17}} \\
TP2（compile 开） & \textcolor{Muted}{0.959} & \textbf{\textcolor{Teal}{0.330}} & \textcolor{Muted}{3.76} & \textbf{\textcolor{Teal}{3.16}} \\
CFG 并行，KV 前缀缓存关 & \textcolor{Muted}{0.826} & \textbf{\textcolor{Teal}{0.293}} & \textcolor{Muted}{3.18} & \textbf{\textcolor{Teal}{2.66}} \\
\midrule
\textcolor{Orange}{Omni two-stage（1386 token 提示），缓存关} & -- & \textbf{\textcolor{Orange}{0.193}} & -- & \textbf{\textcolor{Orange}{1.882}} \\
\textcolor{Orange}{Omni two-stage，Thinker 前缀缓存开} & -- & \textbf{\textcolor{Orange}{0.053}} & -- & \textbf{\textcolor{Orange}{1.744}} \\
\bottomrule
\end{tabularx}

\vspace{4pt}
\begin{itemize}\scriptsize\setlength\itemsep{1pt}
  \item vLLM-Omni two-stage 的图生文请求返回 \code{usage.prompt\_tokens = 1386}（1369 个图像 patch + 17 个文本 token）。M* 的 \code{vit\_patches(512, 512, "default")} 同为 1369。
  \item 早期测试的 M* 图生文使用了 4900 patch（\code{image\_preprocess:"vllm"}），与 vLLM-Omni main 的 1386 token 不等量，标 \textsuperscript{§}，仅供参考。
  \item 灰色 = 早期的 4900 patch 数据；粗体 = 输入对齐后的数据（\code{results/20261122\_i2t\_parity}）。1 token 延迟降低 56--66\%，64 token 降低 4--17\%。
\end{itemize}
\takeaway{对齐后，M* 的 1 token 延迟为 0.085 s（缓存开）、0.293 s（缓存关）；vLLM-Omni 为 0.053 s、0.193 s。两种设置下 vLLM-Omni 都更快。}
\end{frame}

\begin{frame}[t]{图像编辑：同一指令、同一图像 token 量}
\scriptsize\def\codesize{\relax}
\renewcommand{\arraystretch}{1.08}\setlength{\tabcolsep}{3pt}
\begin{tabularx}{\textwidth}{>{\bfseries}p{4.6cm} p{2.5cm} Y}
\toprule
\rowcolor{Mist}
配置 & 图像 token（ViT / VAE） & I2I 中位数（min--max，秒） \\
\midrule
\textcolor{Teal}{M* CFG 并行（compile 开）} & 1369 / 1024 & \textbf{16.79\,\textcolor{Muted}{(16.75--16.81)}} \\
\textcolor{Teal}{M* CFG 并行（compile 关）} & 1369 / 1024 & 18.25\,\textcolor{Muted}{(18.23--18.26)} \\
\textcolor{Teal}{M* TP2（compile 开）} & 1369 / 1024 & 21.21\,\textcolor{Muted}{(21.20--21.21)} \\
\textcolor{Teal}{M* CFG 并行，KV 前缀缓存关} & 1369 / 1024 & 16.82\,\textcolor{Muted}{(16.80--16.83)} \\
\textcolor{Orange}{vLLM-Omni 单 stage TP2（512$^2$ 原生输入）} & 1369 / 1024 & \textbf{28.08\,\textcolor{Muted}{(28.07--28.08)}} \\
\midrule
\multicolumn{3}{l}{\textcolor{Muted}{早期指令（“Add a small cat sitting beside the bagel on the plate.”）的运行}} \\
\textcolor{Muted}{M* CFG 并行，\code{vllm}（4900 patch）} & 4900 / 1024 & 18.98\,\textcolor{Muted}{(18.96--19.00)} \\
\textcolor{Muted}{M* CFG 并行，\code{default}（输出未加入猫）} & 1369 / 1024 & 16.83\,\textcolor{Muted}{(16.81--16.83)} \\
\textcolor{Muted}{vLLM-Omni 单 stage TP2（512$^2$ 原生）} & 1369 / 1024 & 28.07\,\textcolor{Muted}{(28.06--28.07)} \\
\textcolor{Muted}{vLLM-Omni 单 stage TP2（输入放大到 980$^2$）} & 4900 / 3721 & 31.24\,\textcolor{Muted}{(31.24--31.25)} \\
\bottomrule
\end{tabularx}

\vspace{2pt}
\begin{itemize}\scriptsize\setlength\itemsep{0pt}
  \item 上半部分：两边同一指令 \code{Replace the bagel with a croissant.}，ViT 1369 patch，VAE 1024 latent，输出均为可颂。
  \item 下半部分：早期“加猫”指令，仅作记录。M* 默认预处理下未生效，故改用上述指令。
  \item Omni 放大到 980$^2$：VAE 为 3721 latent，不等量，不用于比较。
\end{itemize}
\takeaway{等量后 M* CFG 并行仍快约 40\%：\val{final/mstar-cfgpar/i2i} s，vLLM-Omni 为 \val{final/omni-single-tp2/i2i} s。}
\end{frame}

\begin{frame}[t]{首个请求（含 JIT、编译、图捕获）的延迟，秒}
\scriptsize\def\codesize{\relax}
\renewcommand{\arraystretch}{1.12}\setlength{\tabcolsep}{3pt}
\begin{tabularx}{\textwidth}{>{\bfseries}p{5.2cm} Y Y Y Y}
\toprule
\rowcolor{Mist}
配置 & T2I & I2T 1 token & I2T 64 token & I2I \\
\midrule
M* CFG 并行 · compile 开 · 缓存开 · 无 graph & 17.5 & 8.4 & 3.5 & 34.9 \\
M* CFG 并行 · compile 关（诊断）· 缓存开 & 10.9 & 1.6 & 2.2 & 22.0 \\
M* TP2 · compile 开 · 缓存开 · 无 graph & 24.6 & 16.7 & 4.3 & 39.7 \\
M* CFG 并行 · compile 开 · 缓存关 & -- & 6.7 & 3.6 & 32.3 \\
Omni 单 stage CFG 并行 · 无 compile & 10.9 & -- & -- & -- \\
Omni 单 stage TP2 · 无 graph & 17.7 & -- & -- & 29.5 \\
Omni two-stage · 缓存关 · eager & 18.7 & 0.6 & 1.9 & -- \\
Omni two-stage · 缓存关 · Thinker graph & -- & 0.4 & 1.9 & -- \\
Omni two-stage · 缓存开 · eager & -- & 0.6 & 1.7 & -- \\
\bottomrule
\end{tabularx}

\vspace{4pt}
\begin{itemize}\scriptsize\setlength\itemsep{1pt}
  \item 首请求指服务启动后，该 case 的第一次请求（warmup 1）。第二次 warmup 与稳态已无明显差异。M* 的图生文行来自单独的等量运行；vLLM-Omni 的两行同理。
  \item M* 默认开启 \code{torch.compile}：T2I、I2I、I2T 的首请求均多出约 6--7 s。关闭后首请求更短，但扩散稳态慢 7--9\%。
  \item vLLM-Omni 无 compile，首请求只多出内核懒加载的时间。单 stage CFG 并行的首个 T2I 比稳态多约 1 s。
\end{itemize}
\takeaway{需要冷启动即时响应的场景中，M* 的 compile 会使第一个请求多等约 6--7 s。稳态服务中该开销可以忽略。}
\end{frame}
